\documentclass[a4paper, 11pt]{article}
\usepackage{geometry}
\usepackage{graphicx}
\usepackage{amsmath}
\usepackage{float}
\usepackage{booktabs}

\geometry{margin=1in}

\begin{document}
	
\title{DEEP LEARNING BASED VIRTUAL SENSING METHOD AND SYSTEM FOR LOW-COST AQUACULTURE WATER QUALITY MONITORING}
\maketitle
\section*{Introduction}
The document discusses a novel approach to Water Quality Monitoring (WQM) in aquaculture, utilizing a Virtual Sensing (VS) system integrated with Deep Learning (DL). Specifically, the system employs solubility equations and Deep Learning Recurrent Neural Network (DL-RNN) models within the Dissolved Oxygen Enhanced Virtual Sensing (DLEVS) mechanism.

\section*{Problem Statement}
Effective monitoring of water quality is essential for the management of aquaculture systems, ensuring the well-being and growth of aquatic life. Traditional sensor-based systems are often cost-prohibitive, especially in resource-constrained settings.

\section*{Virtual Sensing Using Solubility Equations}
The VS methodology relies on established solubility equations to estimate water quality parameters, with a primary focus on dissolved oxygen (DO). Mathematical models derived from extensive datasets offer a cost-effective alternative to physical sensors, allowing for flexible estimation irrespective of the water source.

\subsection*{VS Capabilities Across Water Sources}
The VS system, based on solubility equations, is designed to estimate water quality parameters across diverse water sources. Mathematical models enable accurate estimations for scenarios where direct measurements through physical sensors are challenging.

\section*{Deep Learning Enhancement for Accuracy}
To address concerns about the learning mechanism and real-time error correction, the VS system incorporates DL-RNN models through the DLEVS mechanism. This compensates for potential inaccuracies in VS-based DO estimation, particularly in real-time scenarios involving new, untrained water bodies.

\subsection*{DL-RNN Training and Error Prediction}
The DL-RNN model undergoes extensive training on various water quality datasets to predict errors in VS-based DO estimates. This process ensures adaptability and high predictive accuracy across diverse water sources, contributing to the robustness of the DL-RNN.

\subsection*{Real-Time Measurements and DL Enhancement}
In real-time scenarios, salinity and temperature are measured, and the VS system estimates DO using solubility equations. The DL-RNN model predicts errors based on the dynamic characteristics of the specific water source, enhancing adaptability to new water bodies in real time. Subtracting the predicted error from the VS-estimated DO yields an enhanced and more accurate DO estimate reflective of real-time conditions.

\section*{Conclusion}
In summary, this research pioneers an innovative and cost-effective approach to WQM in aquaculture, achieving heightened accuracy through the integration of VS and DL. The system, incorporating solubility equations and DL-RNN models within the DLEVS mechanism, demonstrates exceptional capability in estimating water quality parameters irrespective of the water source.

The integration of DL-RNN models significantly enhances the accuracy of DO estimation by proficiently predicting and compensating for potential errors. Real-time measurements, including salinity and temperature, coupled with the synergy of virtual sensing and DL enhancement, establish a comprehensive and precise approach to water quality estimation.

To the best of the authors’ knowledge, this research marks the first proposal of integrating Virtual Sensing and Deep Learning to reduce the cost of sensor-based WQM systems for aquaculture. The findings contribute to advancing sustainable aquaculture practices by providing a more accessible and accurate method for monitoring water quality parameters.

	
	
	
\end{document}
